\section{The intrinsic classical structure: symplectic and Poisson geometry}

For a configuration manifold \(Q\), the canonical phase space is its cotangent bundle \(T^*Q\). This standard construction is reviewed in Ref.~\cite{RoubtsovDutykh}. In local canonical coordinates \((x^i,p_i)\), the canonical one-form is
\begin{equation}
\theta=p_i\,\dd x^i,
\end{equation}
and the canonical symplectic form is
\begin{equation}
\boxed{
\omega=-\dd\theta=\dd x^i\wedge\dd p_i.
}
\label{eq:symplectic}
\end{equation}
Repeated indices are summed. The inverse of \(\omega\), understood as a bivector, is the canonical Poisson tensor
\begin{equation}
\boxed{
\Pi
=
\frac{\partial}{\partial x^i}
\wedge
\frac{\partial}{\partial p_i}.
}
\label{eq:poisson-bivector}
\end{equation}
It acts on differentials of observables to define the Poisson bracket
\begin{equation}
\boxed{
\pb{f}{g}
=
\Pi(\dd f,\dd g)
=
\frac{\partial f}{\partial x^i}
\frac{\partial g}{\partial p_i}
-
\frac{\partial f}{\partial p_i}
\frac{\partial g}{\partial x^i}.
}
\label{eq:poisson}
\end{equation}
In particular,
\begin{equation}
\boxed{
\pb{x^i}{p_j}=\delta^i{}_j,
\qquad
\pb{x^i}{x^j}=0,
\qquad
\pb{p_i}{p_j}=0.
}
\label{eq:canonical-pb}
\end{equation}

This is the invariant mathematical structure that the cross-gradient ansatz most closely resembles. The wedge product in Eq.~\eqref{eq:poisson-bivector} is antisymmetric, but it is not a three-dimensional vector cross product. It defines a bilinear bracket on observables rather than a universal constant-valued differential operator.

The dimensional bookkeeping is also natural. The symplectic form has dimensions of action, so
\begin{equation}
\frac{\omega}{\hbar}
\end{equation}
is dimensionless. Its inverse \(\Pi\) has dimensions of inverse action, so \(\hbar\Pi\) is dimensionless. This is a legitimate sense in which \(\hbar\) sets the scale of phase-space geometry; it does not imply that every antisymmetric phase-space expression equals a power of \(\hbar\).
